\section{Use a working-set estimate as a prediction}
For a complete \(T\)-by-\(T\)-by-\(T\) multiplication tile, the selected values
occupy \(4T^2\) bytes in each input and \(8T^2\) in the output, or \(16T^2\)
bytes in total. With \(T=4\), this equals the cache's 256-byte capacity. With
\(T=8\), the selected values occupy 1024 bytes. A larger tile contains more
opportunities to reuse a value, but also more values competing to remain resident.

This byte estimate is not a miss-count proof. Subrows can occupy partial cache
lines, the last tiles have different dimensions, and the stack and descriptor
also produce accesses. Placement adds another constraint. For a direct-mapped
sixteen-line cache, a byte address \(x\) selects
\[
 \left\lfloor x/16\right\rfloor\bmod16.
\]
Two live lines with the same set number displace each other. A fully associative
cache allows any resident line to occupy any of its sixteen positions. A two-way
cache has eight sets with two positions per set.

The allocation bases are part of a cache experiment. Moving one matrix by a cache
line preserves the arithmetic but changes its set numbers. Moving a matrix by
only one element can also change which rows straddle line boundaries. Record both
logical shape and actual addresses, together with line size, associativity,
replacement policy and cache reset boundary.

\begin{MARSExercise}{A fit that still conflicts}
For sixteen input columns, neighboring row starts differ by sixty-four bytes.
List the direct-mapped set numbers of the first five row starts when the first
starts in set zero. How many different sets have those five starts used? Explain
why total byte capacity alone does not predict their coexistence.
\end{MARSExercise}
